\chapter{综合练习与术语表}

\section{综合练习一：端到端预测性维护}
某生产线每秒采集温度、压力、振动和电流，维修记录提供故障标签，但标签时间存在延迟。请完成：
\begin{enumerate}[label=\arabic*.]
 \item 定义样本、输入窗口、预测提前量和标签生成规则。
 \item 说明如何按设备和时间划分训练、验证和测试集。
 \item 选择一个序列基线和一个深度模型，并解释选择理由。
 \item 设计针对传感器缺失、工况变化和新设备的三组测试。
 \item 给出误报和漏报成本不同情况下的决策阈值方案。
\end{enumerate}

\section{综合练习二：视觉与文本联合质检}
工人以自然语言记录缺陷，检测相机拍摄产品表面。请比较图像分类、目标检测、图文对比学习和视觉语言生成四种建模方式。讨论标签粒度、文本噪声、缺陷区域定位和模态缺失对结果的影响，并设计至少一个能检验跨产品泛化的划分。

\section{综合练习三：跨工厂联邦迁移}
三个工厂生产相同部件，但传感器型号、采样频率和故障比例不同。请设计一个联邦迁移方案，说明共享参数、个性化参数、客户端采样、隐私保护、掉线处理和最终评价。要求同时报告平均客户端性能和最差客户端性能。

\section{综合练习四：稀缺故障生成}
只有少量真实故障样本，正常样本充足。请比较 VAE、GAN 和扩散模型的适用性。设计生成样本验收流程，至少包含重复性检查、统计分布检查、专家评审和真实测试集下游增益。讨论生成数据可能引入的标签噪声和隐私风险。

\section{综合练习五：少样本新设备适应}
新设备只有每类三个标注样本。请比较线性探测、全量微调、原型网络和 MAML 风格初始化。要求设计任务级训练/测试划分，并说明如何区分快速适应能力与设备身份记忆。

\section{术语表}
\begin{description}
 \item[经验风险] 在有限训练集上计算的平均损失。
 \item[泛化] 模型在未参与训练的数据或任务上的表现。
 \item[校准] 预测概率与实际事件频率的一致程度。
 \item[域偏移] 训练域与部署域的输入或条件分布发生变化。
 \item[负迁移] 使用源域知识后，目标域性能反而下降。
 \item[灾难性遗忘] 学习新任务后旧任务性能显著下降。
 \item[客户端漂移] 联邦学习中本地非 IID 数据使本地更新偏离全局方向。
 \item[流形] 嵌入高维观测空间、局部具有低维结构的集合。
 \item[元学习] 从多个任务中学习如何更快适应新任务。
 \item[参数高效微调] 冻结大部分预训练参数，只训练少量新增或低秩参数。
 \item[模态缺失] 多模态输入中一个或多个数据源不可用。
 \item[后验坍缩] 潜变量模型中近似后验接近先验，潜变量不再携带有效信息。
 \item[安全聚合] 服务器只能得到客户端更新的聚合结果，不能直接读取单个更新。
\end{description}

\section{答题提示}
面对开放题，先画出数据生成和部署流程，再写目标函数。涉及比较时，明确保持哪些因素相同；涉及工业应用时，明确错误成本、时间约束和信息可得性；涉及隐私时，区分数据不上传、更新不可见和形式化差分隐私保证。好的答案不必选择最复杂的模型，但必须让假设、证据和结论彼此对应。